\clearpage

\section*{A. Resumen Ejecutivo}\label{a.-resumen-ejecutivo}
\addcontentsline{toc}{section}{A. Resumen Ejecutivo}

El proyecto consiste en el desarrollo y lanzamiento de una aplicación
web B2C de administración de finanzas personales, instalable en
dispositivos móviles como aplicación web progresiva (PWA), bajo un
modelo de Software como Servicio (SaaS) con opciones freemium. Su
propuesta de valor central reside en la integración profunda con
herramientas de Inteligencia Artificial para simplificar y automatizar
la gestión financiera diaria del usuario. Específicamente, esta
tecnología permite categorizar transacciones de manera automática y
generar recomendaciones personalizadas, basadas en patrones de
comportamiento, para sugerir estrategias financieras. De esta manera, el
producto busca reemplazar métodos manuales o herramientas globales
genéricas mediante una solución intuitiva e inteligente. Además, la
plataforma admite espacios de trabajo compartidos, que permiten que
varias personas colaboren en un mismo seguimiento financiero.

El diferencial competitivo del negocio es su enfoque hiper localizado,
diseñado específicamente para resolver las complejidades únicas de la
realidad económica argentina. La alta inflación crónica y el régimen
bimonetario dificultan enormemente el seguimiento de gastos y el control
patrimonial. Para hacer frente a esta volatilidad, la aplicación realiza
un seguimiento mensual de la inflación y del tipo de cambio del dólar
estadounidense, lo que permite mostrar con precisión la variación en el
valor de las tenencias del usuario. Además, incluye un seguimiento
simple de inversiones, que permite registrar periódicamente sus
ganancias y pérdidas, y un seguimiento de deudas que contempla las
cuotas pendientes de las tarjetas de crédito, los préstamos bancarios y
las deudas con terceros.

Para reducir la fricción en el registro diario de movimientos, la
plataforma opera como un asistente financiero interactivo mediante
modelos de lenguaje. Los usuarios pueden realizar consultas y cargar
gastos de forma omnicanal a través de un chatbot de mensajería
---Telegram durante las etapas de prueba y WhatsApp a partir del
lanzamiento---, utilizando el procesamiento de lenguaje natural para
interpretar mensajes de texto, audios e imágenes. El sistema también
cuenta con ingesta automatizada para extraer datos directamente desde
facturas en formato PDF o tickets de compra. Finalmente, la plataforma
emite alertas predictivas cuando un gasto recurrente aumenta por encima
de la inflación o cuando los gastos estimados para el mes siguiente,
sumados a la deuda y al ajuste por inflación estimado, se acercan o
superan el límite del presupuesto, facilitando la toma de decisiones
financieras informadas en un contexto de inestabilidad constante.

\capitulo{1 Descripción General}

\subsection*{1.1 Descripción básica del
negocio}\label{descripciuxf3n-buxe1sica-del-negocio}
\addcontentsline{toc}{subsection}{1.1 Descripción básica del negocio}

El negocio consiste en el desarrollo y lanzamiento de una aplicación web
de administración de finanzas personales, instalable en dispositivos
móviles como aplicación web progresiva (PWA), operando bajo un modelo de
Software como Servicio (SaaS). La propuesta de valor central reside en
su fuerte integración con capacidades de Inteligencia Artificial (IA),
diseñada para simplificar y automatizar la gestión financiera diaria del
usuario. Específicamente, la IA se utilizará para interpretar mensajes
en lenguaje natural, categorizar automáticamente transacciones y generar
recomendaciones personalizadas basadas en patrones de comportamiento. El
modelo de negocio es B2C (Business-to-Consumer), lo que implica una
estrategia de venta directa a los consumidores finales a través de
suscripciones o un modelo freemium, buscando proporcionar una
herramienta intuitiva que reemplace métodos manuales o soluciones
genéricas.

\subsection*{1.2 Situación actual del
negocio}\label{situaciuxf3n-actual-del-negocio}
\addcontentsline{toc}{subsection}{1.2 Situación actual del negocio}

La situación actual del mercado de gestión de finanzas personales en
Argentina se caracteriza por una marcada complejidad. El seguimiento de
gastos y el control patrimonial son particularmente difíciles debido al
volátil contexto macroeconómico del país. Los desafíos principales
incluyen la alta y crónica inflación, que obliga a los usuarios a
reajustar constantemente sus presupuestos y erosiona el valor real de
los ahorros y los ingresos. Además, el régimen de la doble moneda (Peso
Argentino y Dólar Estadounidense) es un factor crucial: los usuarios
deben manejar activos y deudas en ambas monedas. En los últimos años
surgieron aplicaciones locales que permiten registrar gastos mediante
mensajería e inteligencia artificial y que contemplan ambas monedas; sin
embargo, se concentran en el registro y la categorización de los
movimientos, sin que la inflación aparezca como una variable central
para planificar. Esto deja un espacio para una solución que ofrezca una
visión financiera ajustada a la realidad económica argentina,
facilitando la toma de decisiones informadas a pesar de la
inestabilidad.

En cuanto al estado del emprendimiento, el desarrollo del producto ya se
encuentra en marcha y se cuenta con una base funcional de la aplicación.
Desde la interfaz web es posible gestionar espacios de trabajo, cuentas
y categorías, registrar movimientos y consultar la evolución de los
saldos. Además, un chatbot basado en inteligencia artificial permite
consultar y registrar información a partir de mensajes de texto, audios
y documentos PDF; durante la etapa de desarrollo opera sobre Telegram.
Actualmente, la aplicación se encuentra en etapa de pruebas, utilizada
por el propio emprendedor, y se prevé abrirla a un primer grupo reducido
de usuarios hacia fines de 2026.

\subsection*{1.3 ¿Qué hace único a su
negocio?}\label{quuxe9-hace-uxfanico-a-su-negocio}
\addcontentsline{toc}{subsection}{1.3 ¿Qué hace único a su negocio?}

Si bien existen aplicaciones de finanzas personales adaptadas al mercado
argentino que permiten registrar gastos mediante mensajería e
inteligencia artificial, el registro de movimientos por sí solo ya no
constituye un diferencial. Lo que distingue a Quickexpense es la
incorporación de la inflación como variable central de la gestión
financiera: la plataforma alerta al usuario cuando un gasto recurrente
aumenta por encima del índice de inflación, advierte cuando los gastos
estimados para el mes siguiente ---sumados a la deuda pendiente,
incluidas las cuotas de tarjetas de crédito, y al ajuste por inflación
estimado--- se acercan al límite del presupuesto o lo superan, y muestra
la variación en el valor de las tenencias según la inflación y el tipo
de cambio. De este modo, el usuario no sólo sabe en qué gastó, sino
cuánto le cuesta realmente sostener su nivel de gastos en un contexto de
precios cambiantes.

Adicionalmente, el fuerte enfoque en la implementación de inteligencia
artificial y la facilidad de uso buscan reducir al mínimo la fricción
que experimentan los usuarios en las aplicaciones tradicionales, que
frecuentemente es la debilidad más fuerte de los sistemas de gestión de
finanzas personales. En la encuesta realizada para este proyecto, el
68\% de los participantes manifestó interés en tener más control sobre
sus finanzas, pero sólo si resulta fácil y no le consume mucho tiempo
(\citeproc{ref-luraschi2026encuesta}{Luraschi, 2026}).

\subsection*{1.4 Describa los factores principales que usted considera
harán exitoso su
proyecto}\label{describa-los-factores-principales-que-usted-considera-haruxe1n-exitoso-su-proyecto}
\addcontentsline{toc}{subsection}{1.4 Describa los factores principales
que usted considera harán exitoso su proyecto}

Fundamentalmente, el proyecto depende de la precisión y los costos de
las herramientas de Inteligencia Artificial utilizadas a lo largo de la
aplicación. Si estas no son precisas, los usuarios podrían perder la
confianza en el sistema. Si los costos que representa hacer uso de estas
herramientas de forma significativa son demasiado altos, el modelo de
negocio se vuelve inviable, dado que los usuarios deberían afrontar un
costo de subscripción significativo o ver la función distintiva del
sistema fuertemente limitada.

Asimismo, es fundamental garantizar una buena satisfacción de usuario,
lo que implica que el sistema debe ser confiable, tener una alta
disponibilidad y estar probado correctamente.

\subsection*{1.5 Estrategia: Definir, Misión, Visión y Propósito
Estratégico}\label{estrategia-definir-misiuxf3n-visiuxf3n-y-propuxf3sito-estratuxe9gico}
\addcontentsline{toc}{subsection}{1.5 Estrategia: Definir, Misión,
Visión y Propósito Estratégico}

La misión del proyecto es simplificar la gestión del patrimonio en
Argentina mediante la automatización con Inteligencia Artificial.
Mientras tanto, la visión es convertirse en la herramienta de gestión
financiera estándar en el país.

El propósito bajo el cual se inició el proyecto es devolver al usuario
el control y la tranquilidad sobre su dinero en el contexto de
incertidumbre e inestabilidad crónica que experimenta la región.

La estrategia general del negocio es de enfoque con diferenciación
(\citeproc{ref-porter1980}{Porter, 1980}): concentrarse en un segmento
acotado ---jóvenes y adultos jóvenes afectados por la inflación y la
economía bimonetaria, comenzando por la ciudad de Rosario--- y
diferenciarse mediante herramientas que incorporan la inflación a la
planificación financiera, con una carga de datos de mínima fricción a
través de mensajería e inteligencia artificial.

\subsection*{1.6 Identificación de la oportunidad de
negocio}\label{identificaciuxf3n-de-la-oportunidad-de-negocio}
\addcontentsline{toc}{subsection}{1.6 Identificación de la oportunidad
de negocio}

La oportunidad de negocio surge principalmente de los dos factores que
resultan obstáculos mayores a la hora de utilizar un sistema de gestión
de finanzas personales. El primero es la falta de herramientas que
incorporen el contexto local, y en particular la inflación, a la
planificación financiera. La mayoría de las aplicaciones del ámbito se
diseñaron en contextos económicos diferentes, por lo que al intentar
aplicarlas en el contexto local los usuarios notan que existen varios
aspectos en los que el sistema se ve limitado. En los últimos años
surgieron alternativas locales que contemplan la doble moneda y permiten
registrar gastos mediante mensajería; sin embargo, la inflación sigue
sin ser una variable central de planificación. Ejemplos de esto son la
falta de consideración de la inflación en presupuestos y planes de
ahorro, la ausencia de alertas ante gastos que aumentan por encima del
índice de precios y la dificultad para proyectar el peso de las cuotas y
deudas pendientes en los meses siguientes.

El segundo factor que lleva a la identificación de la oportunidad es la
fricción que existe en la mayoría de aplicaciones, sistemas o métodos
para gestionar gastos, ingresos y movimientos a nivel personal. Estos
frecuentemente involucran el trabajo de cargar manualmente cada acción
que se realiza en el momento o de forma conjunta cada cierto periodo de
tiempo, lo que requiere un esfuerzo significativo por parte del usuario
para asegurarse de que la información quede correctamente registrada.
Versiones simplificadas que trabajan con números generales suelen perder
precisión y no ofrecen al usuario una visión completa de su situación
financiera. La encuesta realizada para este proyecto confirma este
obstáculo: el 35\% de los participantes no registra ninguno de sus
gastos, las herramientas más utilizadas son el cuaderno y las planillas
de cálculo (42\% cada una), y el 81\% dedica menos de una hora por
semana a administrar sus finanzas
(\citeproc{ref-luraschi2026encuesta}{Luraschi, 2026}).

\subsection*{1.7 Capacidades centrales}\label{capacidades-centrales}
\addcontentsline{toc}{subsection}{1.7 Capacidades centrales}

Una de las principales capacidades del emprendedor es el dominio de la
arquitectura para software escalable. Se cuenta con la capacidad para
implementar una arquitectura de monolito de forma tal que permita una
fácil transición a una arquitectura de microservicios si el negocio
demostrara requerirlo en un futuro.

Adicionalmente, se cuenta con conocimiento en la integración profunda de
modelos de inteligencia artificial en aplicaciones, haciendo uso de
herramientas (tool calling) para permitir a un agente realizar
operaciones complejas.

Por otra parte, el emprendedor cuenta con experiencia en el sector
financiero: actualmente trabaja en una empresa que desarrolla tecnología
para un agente de liquidación y compensación (ALyC) y una billetera
virtual, donde participa en el desarrollo de herramientas conectadas al
mercado y de cotización. Esto le aporta una comprensión de las
complejidades macroeconómicas locales y del funcionamiento de los
productos financieros que utilizan los usuarios, que se complementa con
el conocimiento de estrategias de gestión de finanzas personales.

Finalmente, se realizaron consultas a profesionales del área de
contabilidad sobre buenas prácticas para llevar una contabilidad limpia
e íntegra, que se reflejan en el diseño del sistema, basado en el
registro por partida doble.

\subsection*{1.8 Propuesta de valor para el
cliente}\label{propuesta-de-valor-para-el-cliente}
\addcontentsline{toc}{subsection}{1.8 Propuesta de valor para el
cliente}

Un asistente financiero multiplataforma que facilita el registro de
datos mediante IA, ofreciendo herramientas nativas para la economía
argentina como el seguimiento de la inflación y del tipo de cambio, el
seguimiento de deudas y cuotas pendientes, y alertas que advierten
cuando un gasto sube por encima de la inflación o cuando los gastos
estimados se acercan al límite del presupuesto.

\subsection*{1.9 Valores nucleares de la
organización}\label{valores-nucleares-de-la-organizaciuxf3n}
\addcontentsline{toc}{subsection}{1.9 Valores nucleares de la
organización}

La filosofía de la empresa incluye la creencia de que la tecnología debe
representar la menor fricción posible para el ser humano. Esto se
traduce a que el usuario no debe realizar acciones tediosas y
repetitivas, y que el sistema debe ser lo suficientemente inteligente
para traducir lenguaje natural en instrucciones precisas.

Adicionalmente, dada la sensibilidad de la información que se maneja en
el sistema, la seguridad y la privacidad son primordiales para el
proyecto. Esto se ve reflejado en la revisión de arquitecturas desde un
punto de vista de seguridad y el compromiso de utilizar la información
sensible únicamente para proveer las funciones prometidas al usuario.

Finalmente, hay un enfoque en la transparencia al informar al usuario
sobre el riesgo de utilizar herramientas de inteligencia artificial
generativa, así como las implicancias de privacidad que conlleva.

\subsection*{1.10 Enfoque e iniciativas
estratégicas}\label{enfoque-e-iniciativas-estratuxe9gicas}
\addcontentsline{toc}{subsection}{1.10 Enfoque e iniciativas
estratégicas}

La estrategia consiste en ofrecer el servicio mediante múltiples
canales, principalmente en aplicaciones a las que el usuario ya esté
acostumbrado. El canal principal será WhatsApp, la aplicación de
mensajería más utilizada del país
(\citeproc{ref-statista2025redes}{Statista, 2025}), cuya integración
está prevista para el lanzamiento comercial; durante las etapas de
desarrollo y prueba, el asistente opera sobre Telegram. De esta forma se
busca hacer la transición a la aplicación lo más fácil posible.

Al mismo tiempo, se busca diferenciarse tanto de los competidores
extranjeros como de los locales desarrollando y refinando continuamente
herramientas especializadas para el marco económico del país, en
particular aquellas que incorporan la inflación a la planificación
financiera, como el seguimiento de la inflación y del tipo de cambio y
las alertas basadas en esos indicadores.

Por último, el ingreso al mercado será gradual: una versión Alpha con un
grupo reducido de usuarios hacia fines de 2026, una versión Beta a
partir de febrero de 2027 y el lanzamiento previsto para junio de 2027.
Cada etapa estará seguida de un ciclo de retroalimentación y mejora, de
modo de validar la propuesta de valor antes de su comercialización.

\subsection*{1.11 Áreas claves de resultados: identificar y
justificar}\label{uxe1reas-claves-de-resultados-identificar-y-justificar}
\addcontentsline{toc}{subsection}{1.11 Áreas claves de resultados:
identificar y justificar}

Se cuentan con las siguientes áreas clave de resultados:

\begin{itemize}
\tightlist
\item
  \textbf{Costo operativo de API por usuario activo.} Tratándose de un
  SaaS potenciado por IA, el consumo de tokens es el mayor gasto
  variable. Si el costo de las llamadas a modelos de inteligencia
  artificial supera el ingreso promedio de suscripción o el presupuesto
  de marketing para usuarios gratuitos, el negocio quiebra.
\item
  \textbf{Usuarios activos diarios en plataformas de mensajería.} Se
  debe tener en cuenta que la propuesta de valor central es reducir la
  fricción del registro diario de gastos. Si los usuarios no interactúan
  con el asistente de inteligencia artificial mediante las plataformas
  de mensajería casi a diario, significa que la misma no resolvió la
  fricción y que eventualmente abandonarán el producto.
\item
  \textbf{Tasa de conversión de freemium a premium.} Define si las
  herramientas avanzadas generan suficiente valor real como para que el
  usuario argentino esté dispuesto a pagar una suscripción en un mercado
  que atraviesa una recesión.
\item
  \textbf{Latencia de respuesta del asistente.} Los usuarios esperan
  respuestas inmediatas, si el procesamiento de mensajes tarda
  demasiado, la experiencia de usuario se ve fuertemente degradada, lo
  que puede llevarlos a abandonar el producto.
\end{itemize}

\subsection*{1.12 Ingreso al sector: estrategias de
inserción}\label{ingreso-al-sector-estrategias-de-inserciuxf3n}
\addcontentsline{toc}{subsection}{1.12 Ingreso al sector: estrategias de
inserción}

Desde la perspectiva de las estrategias genéricas de Porter
(\citeproc{ref-porter1980}{1980}), el negocio adopta una estrategia de
enfoque con diferenciación. Por un lado, se concentra en un segmento
acotado: jóvenes y adultos jóvenes de 18 a 35 años afectados por la
economía bimonetaria y la alta inflación, comenzando por la ciudad de
Rosario. Por otro lado, no busca competir por precio, sino diferenciarse
mediante herramientas que incorporan la inflación a la planificación
financiera, como las alertas ante gastos recurrentes que aumentan por
encima del índice de precios o la proyección de los gastos del mes
siguiente considerando la deuda pendiente y la inflación estimada.

La estrategia de inserción está completamente enfocada en el nicho de
usuarios afectados por el contexto económico regional. No se busca
competir directamente contra alternativas globales por el mercado
global; en cambio, el negocio se enfoca exclusivamente en usuarios
afectados por la economía bimonetaria, la alta inflación y otras
características de la economía local. Al resolver problemas locales que
las aplicaciones globales ignoran, la plataforma crea una barrera de
entrada natural contra competidores extranjeros. Frente a los
competidores locales, que ya resuelven el registro de gastos mediante
mensajería, la diferenciación se apoya en la profundidad del tratamiento
de la inflación.

Desde la perspectiva de las estrategias competitivas de Kotler \& Keller
(\citeproc{ref-kotler2016}{2016}), la plataforma adopta el rol de
especialista de nicho. Las grandes empresas evitan los mercados con alta
fricción regulatoria o volatilidad extrema porque no escalan fácilmente
a nivel global. Al adoptar esta complejidad, el producto se especializa
en un segmento desatendido.

\capitulo{2 Análisis estratégico}

\subcapitulo{2.1 Análisis de Contexto}

\subsubsection*{2.1.1 Descripción del escenario
local}\label{descripciuxf3n-del-escenario-local}
\addcontentsline{toc}{subsubsection}{2.1.1 Descripción del escenario
local}

El lanzamiento se realizará inicialmente en la ciudad de Rosario,
provincia de Santa Fe. Según el Censo Nacional 2022, el departamento
Rosario cuenta con 1.348.725 habitantes, de los cuales 311.156 tienen
entre 20 y 34 años (\citeproc{ref-correa2023censo}{Correa, 2023}), rango
que constituye el núcleo del público objetivo. Si bien el servicio es
digital y puede utilizarse desde cualquier punto del país, las acciones
de captación se concentrarán en esta ciudad durante la primera etapa.

El escenario se enmarca en el mercado argentino, que, a criterio del
autor, se caracteriza por un alto estrés financiero en la clase media y
en los jóvenes profesionales. La necesidad de control patrimonial es
alta, pero el tiempo y la educación financiera suelen ser bajos:
Argentina se ubicó en el puesto 37 de 39 economías en la medición de
educación financiera realizada por el BCRA y CAF
(\citeproc{ref-bcracaf2018}{Banco Central de la República Argentina \&
CAF - Banco de Desarrollo de América Latina, 2018}), y en la encuesta
realizada para este proyecto los participantes calificaron con un
promedio de 5,45 sobre 10 el control que sienten tener sobre sus
finanzas, mientras que el 81\% declaró dedicarle menos de una hora por
semana (\citeproc{ref-luraschi2026encuesta}{Luraschi, 2026}). Este
escenario resulta propicio para un asistente automatizado.

\subsubsection*{2.1.2 Factores económicos}\label{factores-econuxf3micos}
\addcontentsline{toc}{subsubsection}{2.1.2 Factores económicos}

Argentina atraviesa un escenario de desinflación incompleta combinado
con una marcada debilidad de la actividad económica.

En materia de precios, la inflación de agosto de 2026 fue de 1,7\%
mensual; sin embargo, acumula 21,3\% en los primeros ocho meses del año
y 33,5\% en los últimos doce (\citeproc{ref-indec2026ipc}{Instituto
Nacional de Estadística y Censos, 2026b}). Los precios continúan
variando de forma significativa mes a mes, lo que obliga a los hogares a
revisar sus presupuestos con frecuencia y erosiona el poder adquisitivo
de los ingresos y de los ahorros en pesos.

En cuanto a la actividad, luego de un crecimiento del PBI de 4,4\% en
2025 (\citeproc{ref-lanacion2026pbi}{Manzoni, 2026}), la economía mostró
signos de recesión durante 2026. El PBI del segundo trimestre cayó 0,6\%
desestacionalizado respecto del primero y el consumo privado se contrajo
2,4\% (\citeproc{ref-indec2026pib}{Instituto Nacional de Estadística y
Censos, 2026c}); el Estimador Mensual de Actividad Económica (EMAE) de
julio registró una caída de 2,9\% respecto de junio y de 1,4\%
interanual (\citeproc{ref-indec2026emae}{Instituto Nacional de
Estadística y Censos, 2026a}), y la desocupación ascendió a 7,9\% en el
segundo trimestre (\citeproc{ref-indec2026eph}{Instituto Nacional de
Estadística y Censos, 2026d}). En este contexto, la Fundación Capital
redujo su proyección de crecimiento para 2026 a 1,9\%
(\citeproc{ref-mdz2026emae}{{«Día D para la economía: el Indec revela el
EMAE de julio y se define si Argentina entra en recesión»}, 2026}).

En materia cambiaria, desde abril de 2025 rige un esquema de flotación
entre bandas cuyos límites, desde enero de 2026, se actualizan según el
último dato de inflación disponible
(\citeproc{ref-gimenez2026bandas}{Giménez, 2026}). A comienzos de
octubre de 2026 el dólar mayorista cotizaba a \$1.520, unos 400 pesos
por debajo del techo de la banda (\$1.924,62), y la diferencia entre el
dólar minorista (\$1.540) y el informal (\$1.545) era mínima
(\citeproc{ref-gasalla2026dolar}{Gasalla, 2026}). Aun así, la
coexistencia de varias cotizaciones y la experiencia de devaluaciones
pasadas mantienen al dólar como la principal referencia de valor para
los ahorristas.

Por último, los servicios digitales prestados desde el exterior que se
pagan en pesos tienen un recargo de entre el 51\% y el 56\%: 21\% de
IVA, una percepción de 30\% a cuenta de Ganancias o Bienes Personales y
entre 2\% y 5\% de Ingresos Brutos, según la provincia
(\citeproc{ref-lepiscopo2026iva}{L'Episcopo, 2026}). Esto encarece las
suscripciones de aplicaciones extranjeras cobradas en dólares frente a
alternativas locales cobradas en pesos. En sentido inverso, los
servicios de inteligencia artificial que utiliza la plataforma se
facturan en dólares, lo que expone su estructura de costos a las
variaciones del tipo de cambio.

\subsubsection*{2.1.3 Factores políticos}\label{factores-poluxedticos}
\addcontentsline{toc}{subsubsection}{2.1.3 Factores políticos}

A nivel mundial existe una creciente tendencia a regular el uso de datos
personales y de la inteligencia artificial, y Argentina no es ajena a
ella.

El tratamiento de datos personales en el país se rige por la Ley 25.326
de Protección de los Datos Personales (\citeproc{ref-ley25326}{Ley
25.326 de Protección de los Datos Personales, 2000}), cuya autoridad de
aplicación es la Agencia de Acceso a la Información Pública (AAIP). Dada
la antigüedad de esta norma, en los últimos años se presentaron
proyectos de reforma integral inspirados en el anteproyecto elaborado
por la AAIP y alineados con estándares internacionales como el
Reglamento General de Protección de Datos de la Unión Europea
(\citeproc{ref-szlak2025}{Szlak, 2025}). Durante 2026 ingresaron en la
Cámara de Diputados los expedientes 1751-D-2026
(\citeproc{ref-dpa1751}{Digital Policy Alert, s.~f.}) y 3397-D-2026;
este último incorpora, entre otros aspectos, reglas sobre las decisiones
automatizadas y la elaboración de perfiles, la notificación de
incidentes de seguridad y la figura del delegado de protección de datos
(\citeproc{ref-marval2026datos}{Marval O'Farrell Mairal, 2026}). Ambos
proyectos se encontraban en tratamiento parlamentario al momento de este
relevamiento. En materia de inteligencia artificial, conviven proyectos
que proponen una regulación según el nivel de riesgo de los sistemas
(\citeproc{ref-szlak2025}{Szlak, 2025}) con una iniciativa del Poder
Ejecutivo orientada a promover la actividad
(\citeproc{ref-cronica2026ia}{{«Presentan un proyecto para convertir a
Argentina en un polo global de IA»}, 2026}).

En el ámbito financiero, el Decreto 353/2025 creó el Sistema de Finanzas
Abiertas, que permite a las personas compartir, con su consentimiento
expreso, su información con las entidades del sistema financiero
inscriptas en el Banco Central (\citeproc{ref-decreto353}{Decreto
353/2025, 2025}). El Banco Central, encargado de definir los estándares
del sistema, todavía no había publicado durante 2026 los requisitos
técnicos ni un cronograma de implementación
(\citeproc{ref-fiskil2026}{Fiskil, s.~f.}).

Para el negocio, este marco implica, por un lado, la necesidad de
diseñar el tratamiento de los datos financieros de los usuarios
cumpliendo la normativa vigente y anticipando requisitos más exigentes,
como el consentimiento informado, la transparencia en el uso de
inteligencia artificial y la evaluación de riesgos. Por otro lado, abre
una oportunidad a mediano plazo: si el sistema de finanzas abiertas se
implementa, la plataforma podría importar automáticamente los
movimientos de bancos y billeteras virtuales.

\subsubsection*{2.1.4 Factores
Tecnológicos}\label{factores-tecnoluxf3gicos}
\addcontentsline{toc}{subsubsection}{2.1.4 Factores Tecnológicos}

El contexto tecnológico es favorable para un asistente financiero
conversacional basado en inteligencia artificial.

En cuanto a la conectividad, en octubre de 2025 Argentina contaba con
41,6 millones de usuarios de internet, el 90,6\% de la población, y con
66,6 millones de conexiones móviles activas
(\citeproc{ref-kemp2025digital}{Kemp, 2025}). El teléfono celular es la
tecnología de uso más extendido: lo utiliza el 97,8\% de las personas de
18 a 64 años (\citeproc{ref-indec2025tic}{Instituto Nacional de
Estadística y Censos, 2025}).

La mensajería instantánea ocupa un lugar central en la vida digital del
país: WhatsApp es utilizada mensualmente por el 93,9\% de los usuarios
de internet de 16 a 64 años, y Telegram por el 43,6\%
(\citeproc{ref-statista2025redes}{Statista, 2025}).

Por su parte, la inteligencia artificial generativa ya forma parte de la
vida cotidiana de una porción significativa de la población. Según la
Encuesta Nacional sobre Adopción de IA, realizada entre agosto y
septiembre de 2025, el 45,5\% de los adultos conectados la utiliza con
fines personales, proporción que asciende al 63,9\% en la Generación Z y
al 61,7\% entre quienes completaron estudios universitarios
(\citeproc{ref-utdt2025ia}{Universidad Torcuato Di Tella, 2025}).

Como limitación, los principales bancos y billeteras virtuales del país
no ponen a disposición los datos financieros de sus clientes mediante
APIs abiertas, por lo que la información debe cargarse manualmente o
extraerse de documentos como resúmenes y comprobantes. El Sistema de
Finanzas Abiertas creado en 2025 (\citeproc{ref-decreto353}{Decreto
353/2025, 2025}) podría revertir esta situación a mediano plazo.

\subsubsection*{2.1.5 Descripción del escenario:
escenario--meta}\label{descripciuxf3n-del-escenario-escenariometa}
\addcontentsline{toc}{subsubsection}{2.1.5 Descripción del escenario:
escenario--meta}

El escenario meta es la convergencia de una economía que exige
hiper-vigilancia del dinero diario con la explosión de la Inteligencia
Artificial capaz de realizar esa vigilancia de forma autónoma. El
desarrollo del negocio florecerá mientras persista la inestabilidad
macroeconómica y los LLMs sigan reduciendo sus costos de API.

Los factores analizados afectan al desarrollo del negocio de la
siguiente manera:

\begin{itemize}
\tightlist
\item
  \textbf{Factores económicos.} La inflación, aun en desaceleración,
  sostiene la necesidad de controlar los gastos y de medir su evolución
  frente a los precios, lo que refuerza la propuesta diferencial del
  producto. En cambio, la recesión y la caída del consumo limitan la
  disposición a pagar, lo que justifica un modelo freemium con un plan
  gratuito útil por sí mismo. Los costos de inteligencia artificial,
  facturados en dólares, exponen el margen a las variaciones del tipo de
  cambio, aunque la estabilidad cambiaria actual reduce ese riesgo en el
  corto plazo. Por último, el recargo impositivo sobre las suscripciones
  a aplicaciones extranjeras favorece a una alternativa local cobrada en
  pesos.
\item
  \textbf{Factores políticos.} La normativa de protección de datos
  personales y su posible reforma obligan a diseñar desde el inicio un
  tratamiento transparente y seguro de la información financiera de los
  usuarios. A mediano plazo, la implementación del Sistema de Finanzas
  Abiertas permitiría automatizar la importación de movimientos desde
  bancos y billeteras virtuales.
\item
  \textbf{Factores tecnológicos.} El predominio de WhatsApp define el
  canal principal del lanzamiento, y la alta adopción de inteligencia
  artificial entre los jóvenes reduce la barrera de entrada a un
  asistente conversacional. Mientras los bancos no ofrezcan acceso
  abierto a los datos, la carga mediante lenguaje natural y documentos
  seguirá siendo necesaria, que es precisamente el problema que resuelve
  la plataforma.
\end{itemize}

\subsubsection*{2.1.6 Análisis sectorial. Definición de Oportunidades y
Amenazas del
negocio}\label{anuxe1lisis-sectorial.-definiciuxf3n-de-oportunidades-y-amenazas-del-negocio}
\addcontentsline{toc}{subsubsection}{2.1.6 Análisis sectorial.
Definición de Oportunidades y Amenazas del negocio}

El sector de las aplicaciones de finanzas personales en Argentina se
encuentra en una etapa de crecimiento, con la aparición en los últimos
años de varias propuestas locales que utilizan inteligencia artificial y
aplicaciones de mensajería para simplificar el registro de gastos. Se
trata de un sector todavía fragmentado, sin un actor dominante, en el
que conviven aplicaciones globales, startups locales, las herramientas
incorporadas por las billeteras virtuales y métodos manuales como las
planillas de cálculo.

\textbf{Oportunidades}

\begin{itemize}
\tightlist
\item
  La persistencia de la inflación y la creciente demanda por optimizar
  el dinero frente a la crisis, en un mercado donde las herramientas
  disponibles no incorporan la inflación como variable de planificación.
\item
  La rápida adopción de la inteligencia artificial generativa entre los
  jóvenes, que utilizan cada vez más asistentes conversacionales para
  resolver tareas cotidianas.
\item
  El predominio de WhatsApp como canal digital cotidiano, que permite
  ofrecer el servicio dentro de una aplicación que el usuario ya
  utiliza.
\item
  El recargo impositivo sobre los servicios digitales del exterior
  (\citeproc{ref-lepiscopo2026iva}{L'Episcopo, 2026}), que encarece a
  los competidores globales frente a una alternativa local cobrada en
  pesos.
\item
  A mediano plazo, el Sistema de Finanzas Abiertas
  (\citeproc{ref-decreto353}{Decreto 353/2025, 2025}), que permitiría
  importar automáticamente los movimientos de bancos y billeteras.
\end{itemize}

\textbf{Amenazas}

\begin{itemize}
\tightlist
\item
  La presencia de competidores locales ya establecidos, con propuestas
  de valor similares en la integración de inteligencia artificial y
  WhatsApp y con una base de usuarios consolidada.
\item
  Que gigantes del sector Fintech local decidan desarrollar e integrar
  funciones similares de forma nativa en sus aplicaciones, eliminando
  motivos para utilizar una aplicación dedicada.
\item
  La recesión y la caída del consumo, que reducen la disposición a pagar
  por una suscripción.
\item
  La dependencia de plataformas de terceros: desde enero de 2026, Meta
  prohíbe los chatbots de inteligencia artificial de propósito general
  en la API de WhatsApp Business
  (\citeproc{ref-filgueira2025}{Filgueira, 2025};
  \citeproc{ref-whatsapp-terms}{WhatsApp, s.~f.}); si bien los
  asistentes orientados a tareas específicas siguen permitidos, futuros
  cambios en sus políticas o en sus precios podrían afectar al canal
  principal.
\item
  Los cambios regulatorios en materia de datos personales e inteligencia
  artificial, que podrían imponer nuevas obligaciones de cumplimiento.
\end{itemize}

\subcapitulo{2.2 Análisis de la Competencia}

\subsubsection*{2.2.1 Principales competidores
directos}\label{principales-competidores-directos}
\addcontentsline{toc}{subsubsection}{2.2.1 Principales competidores
directos}

Los principales competidores consisten en servicios globales genéricos,
aplicaciones locales con inteligencia artificial, las herramientas
incorporadas por las billeteras virtuales y sustitutos como planillas de
cálculo.

Dentro de los servicios globales genéricos se encuentran algunos como
Wallet, Spendee e YNAB, que poseen interfaces pulidas, pero fueron
diseñados para otros contextos económicos: no contemplan la inflación
local ni las distintas cotizaciones del dólar
(\citeproc{ref-monthly2026}{Equipo Monthly, 2026};
\citeproc{ref-focusfolio2026}{Focus·Folio, 2026}), e YNAB incluso
recomienda llevar un presupuesto separado para cada moneda
(\citeproc{ref-ynab-monedas}{YNAB, s.~f.-b}). A criterio del autor,
tampoco ofrecen un seguimiento adecuado de las compras en cuotas, una
modalidad de pago muy extendida en el país. Además, se cobran en moneda
extranjera: YNAB, por ejemplo, cuesta USD 14,99 por mes o USD 109 por
año (\citeproc{ref-ynab-pricing}{YNAB, s.~f.-a}), y al pagarse en pesos
el precio aumenta entre un 51\% y un 56\% por impuestos
(\citeproc{ref-lepiscopo2026iva}{L'Episcopo, 2026}).

El grupo más cercano a la propuesta de este proyecto son las
aplicaciones locales que permiten registrar gastos mediante WhatsApp con
inteligencia artificial. La más completa es Gasti, que interpreta
mensajes de texto, audios, imágenes y documentos PDF, maneja pesos,
dólares y criptomonedas, y ofrece gastos recurrentes, presupuestos con
alertas, metas de ahorro, cuentas compartidas e integraciones con
Mercado Pago y otras entidades. Según su sitio, más de 11.000 personas
la utilizan en más de 20 países, y sus planes pagos cuestan USD 6,99 y
USD 9,99 por mes (\citeproc{ref-gasti2026}{Gasti, s.~f.}). Otras
alternativas similares son Piggy, que contempla el dólar blue y MEP, las
cuotas de tarjetas de crédito y la integración con Mercado Pago
(\citeproc{ref-piggy2026}{Equipo Piggy, 2026}), QueWallet
(\citeproc{ref-quewallet}{QueWallet, s.~f.}), Chanchito
(\citeproc{ref-chanchito}{Chanchito, s.~f.}), Finy
(\citeproc{ref-piggy2026}{Equipo Piggy, 2026}) y AI Money
(\citeproc{ref-aimoney}{AI Money, s.~f.}).

Las billeteras virtuales, como Mercado Pago y Naranja X, incorporan
herramientas básicas de organización, como la separación del dinero en
compartimentos (los ``frascos'' de Naranja X), aunque limitadas a los
movimientos de su propia plataforma.

Finalmente, las alternativas como planillas de cálculo o libretas de
papel son altamente personalizables, pero tienen un nivel extremo de
fricción para la carga manual de datos. Aun así, son las herramientas
más utilizadas por los participantes de la encuesta realizada para este
proyecto, con un 42\% cada una
(\citeproc{ref-luraschi2026encuesta}{Luraschi, 2026}).

{\def\LTcaptype{none} % do not increment counter
\begin{longtable}[]{@{}
  >{\raggedright\arraybackslash}p{(\linewidth - 8\tabcolsep) * \real{0.2024}}
  >{\raggedright\arraybackslash}p{(\linewidth - 8\tabcolsep) * \real{0.1429}}
  >{\raggedright\arraybackslash}p{(\linewidth - 8\tabcolsep) * \real{0.1429}}
  >{\raggedright\arraybackslash}p{(\linewidth - 8\tabcolsep) * \real{0.1905}}
  >{\raggedright\arraybackslash}p{(\linewidth - 8\tabcolsep) * \real{0.2857}}@{}}
\toprule\noalign{}
\begin{minipage}[b]{\linewidth}\raggedright
Competidor
\end{minipage} & \begin{minipage}[b]{\linewidth}\raggedright
Tipo
\end{minipage} & \begin{minipage}[b]{\linewidth}\raggedright
Canales
\end{minipage} & \begin{minipage}[b]{\linewidth}\raggedright
Inflación como variable de planificación
\end{minipage} & \begin{minipage}[b]{\linewidth}\raggedright
Precio
\end{minipage} \\
\midrule\noalign{}
\endhead
\bottomrule\noalign{}
\endlastfoot
YNAB & Global & App / web & No & USD 14,99/mes o USD 109/año, más
impuestos \\
Spendee & Global & App & No & Suscripción en moneda extranjera (varía
según la región), más impuestos \\
Wallet (BudgetBakers) & Global & App / web & No & Suscripción en moneda
extranjera (varía según la región), más impuestos \\
Gasti & Local & WhatsApp, iOS & No & Gratis; Pro USD 6,99/mes; Premium
USD 9,99/mes \\
Piggy & Local & App + WhatsApp & No & Gratis y planes pagos \\
Mercado Pago / Naranja X & Billetera virtual & App propia & No &
Gratis \\
Planillas de cálculo / cuaderno & Sustituto &
\begin{minipage}[t]{\linewidth}\raggedright
\begin{center}\rule{0.5\linewidth}{0.5pt}\end{center}
\end{minipage} & Manual & Gratis \\
\end{longtable}
}

\subsubsection*{2.2.2 Análisis de cadena de
valor}\label{anuxe1lisis-de-cadena-de-valor}
\addcontentsline{toc}{subsubsection}{2.2.2 Análisis de cadena de valor}

El valor de la plataforma no se genera mediante el almacenamiento de los
datos, si no con la traducción e interpretación. La cadena de valor
central reside en la ingesta automatizada de datos no estructurados a
través de IA, el procesamiento y enriquecimiento local de datos, y la
devolución de reportes proactivos en tiempo real al usuario.

\subsubsection*{2.2.3 Definición de Factores Críticos de Éxito
(FCE)}\label{definiciuxf3n-de-factores-cruxedticos-de-uxe9xito-fce}
\addcontentsline{toc}{subsubsection}{2.2.3 Definición de Factores
Críticos de Éxito (FCE)}

Para lograr que el usuario opte por usar el servicio propio se debe
garantizar la precisión de las herramientas de IA; las herramientas que
se le ofrecen al LLM no pueden fallar al registrar montos, un error
pequeño podría romper la confianza del usuario.

Adicionalmente se debe mantener la disponibilidad multicanal de baja
latencia, asegurándose de que el chatbot responda de forma rápida y que
la infraestructura soporte el uso del sistema tanto via aplicaciones de
mensajería como la propia.

Finalmente, debe haber una actualización dinámica del contexto
macroeconómico. Las APIs que alimentan los tipos de cambio y los índices
de inflación deben ser precisas y estar siempre en línea.

\subsubsection*{2.2.4 Fortalezas y debilidades del
negocio}\label{fortalezas-y-debilidades-del-negocio}
\addcontentsline{toc}{subsubsection}{2.2.4 Fortalezas y debilidades del
negocio}

Las principales fortalezas del negocio son el diseño hiperlocalizado de
la lógica financiera, que incorpora la inflación a la planificación de
gastos, presupuestos y deudas; la reducción a casi cero de la fricción
de uso gracias al procesamiento de lenguaje natural; una base contable
rigurosa, sustentada en el registro por partida doble; y la experiencia
del emprendedor en el desarrollo de tecnología para el mercado
financiero local. A esto se suma una estructura de costos fijos muy baja
en la etapa inicial.

La debilidad más significativa es la dependencia estructural de
proveedores de terceros de infraestructura de IA (como Anthropic y
OpenAI), exponiendo el margen de ganancia a fluctuaciones en los costos
de tokens utilizados para responder a solicitudes de usuarios. A esto se
suman los recursos limitados de un equipo unipersonal y el ingreso
tardío a un mercado en el que ya operan competidores locales con una
base de usuarios establecida.

\capitulo{3 Análisis FODA}

\subsection*{3.1 Cuadro FODA}\label{cuadro-foda}
\addcontentsline{toc}{subsection}{3.1 Cuadro FODA}

{\def\LTcaptype{none} % do not increment counter
\begin{longtable}[]{@{}
  >{\raggedright\arraybackslash}p{(\linewidth - 2\tabcolsep) * \real{0.5000}}
  >{\raggedright\arraybackslash}p{(\linewidth - 2\tabcolsep) * \real{0.5000}}@{}}
\toprule\noalign{}
\endhead
\bottomrule\noalign{}
\endlastfoot
\begin{minipage}[t]{\linewidth}\raggedright
\textbf{Fortalezas} (internas)

\begin{itemize}
\tightlist
\item
  \textbf{Hiper-localización:} lógica financiera adaptada nativamente a
  la economía argentina, que incorpora la inflación a la planificación
  (alertas ante gastos que superan la inflación, proyección del
  presupuesto con inflación estimada, seguimiento de deudas y cuotas,
  seguimiento de la inflación y del tipo de cambio).
\item
  \textbf{Fricción mínima vía IA:} automatización del registro de gastos
  a través de lenguaje natural (texto, audios, imágenes y documentos)
  mediante un chatbot de mensajería.
\item
  \textbf{Conocimiento técnico cruzado:} dominio tanto del desarrollo de
  software (integración de IA, arquitecturas escalables) como del
  ecosistema financiero local, a partir de la experiencia del
  emprendedor en tecnología para un agente de liquidación y compensación
  y una billetera virtual.
\item
  \textbf{Base contable rigurosa:} registro por partida doble, diseñado
  siguiendo buenas prácticas consultadas con profesionales de la
  contabilidad.
\item
  \textbf{Costos fijos bajos:} la operación inicial requiere alrededor
  de USD 80 mensuales.
\end{itemize}
\end{minipage} & \begin{minipage}[t]{\linewidth}\raggedright
\textbf{Debilidades} (internas)

\begin{itemize}
\tightlist
\item
  \textbf{Estructura de costos dolarizada:} los costos variables
  dependen del consumo de tokens de APIs de terceros (como OpenAI o
  Anthropic), lo cual presenta un riesgo frente a devaluaciones locales.
\item
  \textbf{Recursos limitados:} el desarrollo, el mantenimiento técnico y
  el soporte dependen de un equipo unipersonal, lo que limita la
  velocidad de escalabilidad técnica.
\item
  \textbf{Ingreso tardío:} el producto todavía no fue lanzado, mientras
  que existen competidores locales con una base de usuarios establecida.
\item
  \textbf{Canal principal en desarrollo:} la integración con WhatsApp
  todavía no está implementada.
\end{itemize}
\end{minipage} \\
\begin{minipage}[t]{\linewidth}\raggedright
\textbf{Oportunidades} (externas)

\begin{itemize}
\tightlist
\item
  \textbf{Adopción temprana de IA:} un público joven que está
  normalizando rápidamente el uso de asistentes conversacionales para
  resolver tareas cotidianas (63,9\% de uso en la Generación Z)
  (\citeproc{ref-utdt2025ia}{Universidad Torcuato Di Tella, 2025}).
\item
  \textbf{Necesidad de control financiero:} la inflación (33,5\%
  interanual en agosto de 2026) (\citeproc{ref-indec2026ipc}{Instituto
  Nacional de Estadística y Censos, 2026b}) obliga a los usuarios a
  llevar un control estricto de su patrimonio para evitar la pérdida de
  poder adquisitivo.
\item
  \textbf{Predominio de WhatsApp:} la utiliza el 93,9\% de los usuarios
  de internet del país (\citeproc{ref-statista2025redes}{Statista,
  2025}).
\item
  \textbf{Recargo impositivo a las aplicaciones extranjeras:} entre 51\%
  y 56\% al pagarlas en pesos
  (\citeproc{ref-lepiscopo2026iva}{L'Episcopo, 2026}).
\item
  \textbf{Finanzas abiertas:} la futura implementación del sistema
  permitiría importar movimientos automáticamente.
\end{itemize}
\end{minipage} & \begin{minipage}[t]{\linewidth}\raggedright
\textbf{Amenazas} (externas)

\begin{itemize}
\tightlist
\item
  \textbf{Competencia directa establecida:} startups locales con
  propuestas de valor similares en la integración de IA y WhatsApp (por
  ejemplo, Gasti, que declara más de 11.000 personas usuarias
  (\citeproc{ref-gasti2026}{Gasti, s.~f.}), o Piggy).
\item
  \textbf{Evolución de las billeteras virtuales:} el riesgo de que
  actores masivos (como Naranja X con sus ``frascos'' o Mercado Pago)
  desarrollen e integren herramientas de presupuestación más complejas o
  asistentes nativos en sus plataformas.
\item
  \textbf{Recesión:} la caída del consumo reduce la disposición a pagar
  suscripciones.
\item
  \textbf{Dependencia de plataformas de terceros:} cambios en las
  políticas o precios de WhatsApp Business o de los proveedores de IA.
\item
  \textbf{Cambios regulatorios} en materia de datos personales e
  inteligencia artificial.
\end{itemize}
\end{minipage} \\
\end{longtable}
}

\subsection*{3.2 Análisis de las Fortalezas, Oportunidades, Debilidades
y
Amenazas}\label{anuxe1lisis-de-las-fortalezas-oportunidades-debilidades-y-amenazas}
\addcontentsline{toc}{subsection}{3.2 Análisis de las Fortalezas,
Oportunidades, Debilidades y Amenazas}

El cruce de estas variables revela una estrategia clara de mitigación y
ataque, que puede organizarse en cuatro tipos de estrategias.

\textbf{Estrategias ofensivas (fortalezas y oportunidades).} La
hiper-localización y la fricción mínima permiten aprovechar la necesidad
de control financiero generada por la inflación y la adopción de la
inteligencia artificial entre los jóvenes. La integración con WhatsApp
para el lanzamiento permite llegar al público en el canal que ya utiliza
a diario, y el cobro en pesos posiciona al producto frente a las
aplicaciones extranjeras, encarecidas por los impuestos.

\textbf{Estrategias defensivas (fortalezas y amenazas).} Frente a la
amenaza de las billeteras virtuales establecidas, que ofrecen
herramientas básicas como la separación del dinero en ``frascos'', el
proyecto opone su principal fortaleza: la ``fricción cero''. Las
billeteras exigen procesos manuales y se limitan a sus propios
movimientos, mientras que este sistema actúa como un asistente autónomo
y proactivo, creando una barrera defensiva basada en la extrema
comodidad del usuario. Frente a los competidores locales que ya
registran gastos por mensajería, la defensa es la profundidad del
tratamiento de la inflación.

\textbf{Estrategias adaptativas (debilidades y oportunidades).} La
principal debilidad del proyecto, los costos dolarizados de las APIs de
IA, se contrarresta mediante la implementación de un modelo Freemium
estricto, que limita las interacciones diarias para usuarios gratuitos y
reserva el acceso ilimitado o a modelos de razonamiento más complejos
exclusivamente para los usuarios de pago. A futuro, el conocimiento
técnico permite contemplar una migración hacia modelos open-source
locales para reducir dependencias. Las limitaciones de un equipo
unipersonal se compensan con un lanzamiento gradual (Alpha, Beta y
lanzamiento), que acota la cantidad de usuarios en cada etapa.

\textbf{Estrategias de supervivencia (debilidades y amenazas).} El
ingreso tardío frente a competidores establecidos y la recesión se
enfrentan con una inserción acotada a un nicho geográfico (Rosario) y
con costos fijos mínimos, que permiten sostener la operación sin
financiamiento externo mientras se valida la propuesta.

\subsection*{3.3 Conclusión: Atractivo de la Industria, Fortalezas del
Negocio}\label{conclusiuxf3n-atractivo-de-la-industria-fortalezas-del-negocio}
\addcontentsline{toc}{subsection}{3.3 Conclusión: Atractivo de la
Industria, Fortalezas del Negocio}

La industria de las Finanzas Personales (PFM) en Argentina presenta un
alto atractivo debido a la ineficacia de las soluciones tradicionales
frente a la volatilidad macroeconómica y a la persistencia de la
inflación. Al mismo tiempo, se trata de un sector en crecimiento en el
que ya operan propuestas locales basadas en inteligencia artificial y
mensajería, por lo que el atractivo no reside en la novedad del canal,
sino en la capacidad de resolver mejor las necesidades específicas del
contexto argentino.

La fortaleza central del negocio es su capacidad para resolver la
``paradoja del usuario moderno'': la necesidad de llevar un control
exhaustivo y casi obsesivo del dinero, pero con una nula disposición a
invertir tiempo en la carga manual de datos. Al posicionar a la
Inteligencia Artificial como el puente entre el control total y el
esfuerzo nulo, y al incorporar la inflación como variable central de
planificación, el proyecto posee una ventaja competitiva altamente
valorada en un mercado inflacionario.

\capitulo{4 Segmentación}

\subsection*{4.1 Segmentación de consumidores y/o
Negocios}\label{segmentaciuxf3n-de-consumidores-yo-negocios}
\addcontentsline{toc}{subsection}{4.1 Segmentación de consumidores y/o
Negocios}

El modelo de negocio es estrictamente B2C (Business-to-Consumer),
orientado a usuarios individuales finales. Si bien el foco está puesto
en la gestión de las finanzas personales de cada usuario, la plataforma
admite espacios de trabajo compartidos, lo que permite que varias
personas colaboren en un mismo seguimiento financiero.

\subsection*{4.2 Identificación de grupos diferenciados de
consumidores}\label{identificaciuxf3n-de-grupos-diferenciados-de-consumidores}
\addcontentsline{toc}{subsection}{4.2 Identificación de grupos
diferenciados de consumidores}

\textbf{Factores Demográficos y Socioculturales:} El público objetivo
principal está compuesto por jóvenes y adultos jóvenes (Generación Z y
Millennials más jóvenes, de 18 a 35 años). Dentro de este rango se
distinguen dos grupos: quienes se encuentran en la etapa de transición
hacia la independencia económica o gestionando sus primeros ingresos
estables, y profesionales ya establecidos en su actividad, con una
estructura financiera más desarrollada que requiere un seguimiento más
exhaustivo.

\textbf{Ubicación geográfica:} en la primera etapa, el público se
concentra en la ciudad de Rosario, cuyo departamento cuenta con 311.156
habitantes de entre 20 y 34 años según el Censo 2022
(\citeproc{ref-correa2023censo}{Correa, 2023}).

\textbf{Nivel educativo y socioeconómico:} la adopción de inteligencia
artificial es mayor entre quienes tienen educación superior (61,7\%
entre universitarios completos, frente a 29,5\% entre quienes no
completaron el secundario) (\citeproc{ref-utdt2025ia}{Universidad
Torcuato Di Tella, 2025}), por lo que el público con estudios terciarios
o universitarios, completos o en curso, es el más receptivo a la
propuesta. En la encuesta realizada para este proyecto, el 71\% de los
participantes cursa estudios universitarios, el 45\% no trabaja y el
52\% vive con su familia sin afrontar gastos en necesidades básicas
(\citeproc{ref-luraschi2026encuesta}{Luraschi, 2026}), lo que indica que
el primer grupo tiene una capacidad de pago limitada y gastos de menor
complejidad.

\textbf{Factores de Comportamiento Tecnológico:} Son individuos nativos
digitales que ya han incorporado la Inteligencia Artificial generativa a
su vida académica o laboral (el 63,9\% de la Generación Z la utiliza con
fines personales (\citeproc{ref-utdt2025ia}{Universidad Torcuato Di
Tella, 2025})), y que utilizan plataformas de mensajería, principalmente
WhatsApp, como su centro operativo digital diario.

\subsection*{4.3 ¿Quiénes son los potenciales usuarios/ compradores del
negocio?}\label{quiuxe9nes-son-los-potenciales-usuarios-compradores-del-negocio}
\addcontentsline{toc}{subsection}{4.3 ¿Quiénes son los potenciales
usuarios/ compradores del negocio?}

El usuario objetivo principal es aquel joven que recién se introduce en
el mundo de la gestión patrimonial. Al no tener un sistema predefinido
(como complejas planillas de Excel heredadas o años de uso de otra
plataforma), representan un mercado ``virgen'' más fácil de captar.

Un segundo grupo está conformado por profesionales ya establecidos en su
actividad, con una estructura financiera más desarrollada. Para este
grupo, el valor de la plataforma reside en reducir el esfuerzo que
implica mantener un seguimiento exhaustivo de su situación financiera.

El comprador (aquel que convierte del modelo gratuito al modelo Premium)
se define por la intensidad de uso: es el usuario que encuentra un valor
indispensable en el asistente conversacional, alcanza rápidamente su
límite de interacciones gratuitas diarias y está dispuesto a pagar una
suscripción para desbloquear el uso ilimitado y el razonamiento avanzado
sin fricciones. En la encuesta realizada para este proyecto, el 52\% de
los participantes manifestó estar dispuesto a pagar una suscripción
mensual por una herramienta que le ahorrara el trabajo manual de
registro (\citeproc{ref-luraschi2026encuesta}{Luraschi, 2026}),
proporción que asciende al 59\% entre quienes trabajan y al 70\% entre
quienes están dispuestos a esforzarse por mejorar el control de sus
finanzas. Se trata de una intención declarada, que suele ser superior a
la conversión efectiva.

\subsection*{4.4 Pautas de comportamiento esperado de cada
segmento}\label{pautas-de-comportamiento-esperado-de-cada-segmento}
\addcontentsline{toc}{subsection}{4.4 Pautas de comportamiento esperado
de cada segmento}

El perfil no se segmenta por sexo ni por un nivel de ingreso específico.
Se proyecta un perfil de usuario con una necesidad de ``micro-control
obsesivo'', que desea saber exactamente adónde va cada centavo, pero
cuya tolerancia a la fricción de uso es casi nula. El comportamiento
esperado consiste en sesiones de uso extremadamente cortas
(micro-interacciones) múltiples veces al día. El usuario no abrirá la
aplicación para navegar por menús; en su lugar, enviará un audio rápido
o una fotografía de un ticket al chatbot inmediatamente después de
realizar un consumo. El esfuerzo cognitivo del usuario se reduce a cero,
trasladando la carga de la categorización, el cálculo de cuotas y el
impacto inflacionario directamente a la inteligencia artificial.

La encuesta realizada para este proyecto respalda este perfil
(\citeproc{ref-luraschi2026encuesta}{Luraschi, 2026}). El 81\% de los
participantes dedica menos de una hora por semana a sus finanzas y el
68\% querría tener más control, pero sólo si resulta fácil y no le
consume tiempo. La principal razón de uso esperada es saber con
exactitud cuánto se gasta por mes y en qué (81\%), seguida por conocer
la rentabilidad de las inversiones (48\%), gestionar el ahorro para
alcanzar objetivos (45\%) y armar un presupuesto con mayor confianza
(42\%). La mayoría utiliza entre dos y tres cuentas por semana (65\%), y
la carga de gastos mediante WhatsApp resulta igual o más útil que la
anotación manual para el 81\%.

Se espera un comportamiento diferenciado entre los dos grupos del
segmento. En el grupo más joven predominan los gastos de menor
complejidad: el 61\% no utiliza tarjeta de crédito y la mediana de uso
de moneda extranjera es de 1 sobre 10, por lo que su uso se centrará en
el registro y la categorización de gastos. En el grupo de profesionales
establecidos se espera un uso más intensivo de las funciones
diferenciales, como el seguimiento de deudas y cuotas, el presupuesto
ajustado por inflación y el valor de las tenencias en distintas monedas.

\capitulo{5 Plan de Acción}

\subsection*{5.1 Programas generales de
acción}\label{programas-generales-de-acciuxf3n}
\addcontentsline{toc}{subsection}{5.1 Programas generales de acción}

La estrategia formulada en los capítulos anteriores se lleva a la
práctica mediante seis programas generales de acción, que se desarrollan
de forma gradual a lo largo de las etapas Alpha, Beta y lanzamiento.

\begin{enumerate}
\def\labelenumi{\arabic{enumi}.}
\tightlist
\item
  \textbf{Programa de desarrollo del producto.} Completar el producto
  mínimo viable y evolucionarlo en cada etapa, priorizando el registro
  de movimientos mediante lenguaje natural y las funciones diferenciales
  que incorporan la inflación a la planificación financiera.
\item
  \textbf{Programa de validación con usuarios.} Abrir el producto a
  grupos acotados de usuarios: menos de 20 en la versión Alpha y entre
  20 y 50 en la versión Beta. Cada etapa se cierra con un ciclo de
  retroalimentación (entrevistas y análisis de uso) que define los
  refinamientos y las nuevas funciones de la siguiente. La convocatoria
  inicial parte de los participantes de la encuesta que manifestaron
  interés en el proyecto (\citeproc{ref-luraschi2026encuesta}{Luraschi,
  2026}).
\item
  \textbf{Programa de integración de canales.} Incorporar WhatsApp como
  canal principal antes del lanzamiento comercial, ajustando el
  asistente a las políticas de la plataforma. Durante las etapas de
  prueba el asistente opera sobre Telegram.
\item
  \textbf{Programa de lanzamiento y captación.} Lanzar el producto en
  junio de 2027, con acciones de comunicación concentradas en la ciudad
  de Rosario.
\item
  \textbf{Programa de monetización.} Activar el modelo freemium en el
  lanzamiento: un plan gratuito con límites de uso y un plan pago con
  acceso ilimitado y funciones avanzadas.
\item
  \textbf{Programa de operación y control de costos.} Sostener la
  infraestructura sobre un servidor dedicado propio y monitorear el
  costo de inteligencia artificial por usuario activo, principal costo
  variable del negocio.
\end{enumerate}

El proyecto se financia íntegramente con capital propio. La inversión
es, principalmente, el tiempo de desarrollo del emprendedor; los costos
monetarios estimados para la etapa inicial son de alrededor de USD 30
mensuales por el servidor dedicado y USD 50 mensuales en consumo de
modelos de inteligencia artificial.

\subsection*{5.2 Programas específicos de
acción}\label{programas-especuxedficos-de-acciuxf3n}
\addcontentsline{toc}{subsection}{5.2 Programas específicos de acción}

El siguiente cronograma detalla las acciones previstas desde la etapa
actual de desarrollo hasta el cierre del primer año de operación
comercial.

{\def\LTcaptype{none} % do not increment counter
\begin{longtable}[]{@{}
  >{\raggedright\arraybackslash}p{(\linewidth - 6\tabcolsep) * \real{0.2439}}
  >{\raggedright\arraybackslash}p{(\linewidth - 6\tabcolsep) * \real{0.1463}}
  >{\raggedright\arraybackslash}p{(\linewidth - 6\tabcolsep) * \real{0.3659}}
  >{\raggedright\arraybackslash}p{(\linewidth - 6\tabcolsep) * \real{0.2195}}@{}}
\toprule\noalign{}
\begin{minipage}[b]{\linewidth}\raggedright
Etapa
\end{minipage} & \begin{minipage}[b]{\linewidth}\raggedright
Período
\end{minipage} & \begin{minipage}[b]{\linewidth}\raggedright
Acciones principales
\end{minipage} & \begin{minipage}[b]{\linewidth}\raggedright
Hito / entregable
\end{minipage} \\
\midrule\noalign{}
\endhead
\bottomrule\noalign{}
\endlastfoot
Desarrollo del MVP & Octubre a diciembre de 2026 & Completar las
funciones de la versión Alpha (detalladas más abajo). Preparar el
entorno de producción. Convocar a los usuarios de la versión Alpha. &
MVP funcional \\
Alpha & Fines de 2026 & Apertura a menos de 20 usuarios, convocados
entre los participantes de la encuesta que dejaron un contacto. &
Primeros usuarios externos \\
Retroalimentación 1 & Enero de 2027 & Entrevistas y análisis de uso.
Corrección de errores, refinamiento de funciones y desarrollo de nuevas
funciones. & Versión Beta \\
Beta & Febrero de 2027 & Apertura a un grupo de 20 a 50 usuarios. &
Versión pre-lanzamiento \\
Retroalimentación 2 y pulido & Marzo a mayo de 2027 & Refinamiento y
pulido. Integración del canal WhatsApp y desarrollo de las funciones
previstas para el lanzamiento. Implementación del plan pago y
preparación de la comunicación del lanzamiento. & Versión de
lanzamiento \\
Lanzamiento & Junio de 2027 & Lanzamiento comercial en Rosario y
activación del modelo freemium. & Producto disponible al público \\
Primer año de operación & Junio de 2027 a mayo de 2028 & Acciones de
captación en Rosario, seguimiento trimestral de las metas comerciales
(6.1) y mejora continua del producto. & Al menos 100 usuarios
registrados \\
\end{longtable}
}

El cronograma se resume en el siguiente diagrama de Gantt, en el que los
rombos señalan los hitos de cada etapa:

\begin{center}
\begin{ganttchart}[
  hgrid, vgrid,
  time slot format=isodate-yearmonth,
  time slot unit=month,
  x unit=0.5cm, y unit chart=0.6cm, y unit title=0.55cm,
  title height=1,
  title label font=\scriptsize,
  bar label font=\small,
  milestone label font=\small\itshape,
  bar/.append style={fill=gray!45},
  milestone/.append style={fill=black},
  bar height=0.55
]{2026-10}{2028-05}
\gantttitle{2026}{3} \gantttitle{2027}{12} \gantttitle{2028}{5} \\
\gantttitle{O}{1} \gantttitle{N}{1} \gantttitle{D}{1} \gantttitle{E}{1} \gantttitle{F}{1} \gantttitle{M}{1} \gantttitle{A}{1} \gantttitle{M}{1} \gantttitle{J}{1} \gantttitle{J}{1} \gantttitle{A}{1} \gantttitle{S}{1} \gantttitle{O}{1} \gantttitle{N}{1} \gantttitle{D}{1} \gantttitle{E}{1} \gantttitle{F}{1} \gantttitle{M}{1} \gantttitle{A}{1} \gantttitle{M}{1} \\
\ganttbar{Desarrollo del MVP}{2026-10}{2026-12} \\
\ganttmilestone{Alpha}{2026-12} \\
\ganttbar{Retroalimentación 1}{2027-01}{2027-01} \\
\ganttbar{Beta}{2027-02}{2027-02} \\
\ganttbar{Retroalimentación 2 y pulido}{2027-03}{2027-05} \\
\ganttmilestone{Lanzamiento}{2027-06} \\
\ganttbar{Primer año de operación}{2027-06}{2028-05}
\end{ganttchart}
\end{center}

Las funciones del producto se incorporan de forma escalonada según la
etapa:

{\def\LTcaptype{none} % do not increment counter
\begin{longtable}[]{@{}
  >{\raggedright\arraybackslash}p{(\linewidth - 2\tabcolsep) * \real{0.2083}}
  >{\raggedright\arraybackslash}p{(\linewidth - 2\tabcolsep) * \real{0.7778}}@{}}
\toprule\noalign{}
\begin{minipage}[b]{\linewidth}\raggedright
Etapa
\end{minipage} & \begin{minipage}[b]{\linewidth}\raggedright
Funciones
\end{minipage} \\
\midrule\noalign{}
\endhead
\bottomrule\noalign{}
\endlastfoot
Alpha (fines de 2026) & \textbf{Ya disponibles:} cuentas, movimientos y
categorías; chatbot con mensajes de texto, audios y documentos PDF;
espacios de trabajo compartidos con roles; evolución de saldos; cuentas
en distintas monedas. \textbf{A incorporar:} interpretación de fotos de
tickets; categorización automática de gastos; invitación de miembros por
email; seguimiento mensual de la inflación; seguimiento del tipo de
cambio del dólar; reportes mensuales de gastos; gastos recurrentes;
presupuesto mensual asistido por inteligencia artificial; metas de
ahorro; deuda de tarjetas de crédito y cuotas pendientes; otras deudas
(con familiares o amigos); instalación en el celular como aplicación web
progresiva (PWA). \\
Lanzamiento (junio de 2027) & Integración con WhatsApp; alerta ante
gastos recurrentes que aumentan por encima de la inflación; deuda de
préstamos bancarios; seguimiento simple de inversiones; carga masiva de
movimientos desde resúmenes de tarjeta o de banco; plan pago y cobro de
la suscripción. \\
Posteriores al lanzamiento & Alerta ante gastos estimados que se acercan
al límite del presupuesto o lo superan; recomendaciones
personalizadas. \\
\end{longtable}
}

El cronograma se extenderá hasta alcanzar el punto de equilibrio
operativo, cuya fecha se determinará en el plan financiero.

\capitulo{6 Plan de Marketing}

\subsection*{6.1 Objetivos}\label{objetivos}
\addcontentsline{toc}{subsection}{6.1 Objetivos}

Los objetivos comerciales se plantean para el primer año de operación,
comprendido entre el lanzamiento (junio de 2027) y mayo de 2028, y se
apoyan en la investigación de mercado primaria realizada para este
proyecto y en indicadores de referencia de la industria.

\textbf{Investigación de mercado.} Entre el 19 y el 20 de septiembre de
2026 se realizó una encuesta en línea a 31 personas de entre 19 y 27
años, en su mayoría estudiantes universitarios, seleccionadas por
conveniencia, por lo que sus resultados son orientativos
(\citeproc{ref-luraschi2026encuesta}{Luraschi, 2026}). Los principales
hallazgos son:

\begin{itemize}
\tightlist
\item
  El 35\% no registra ninguno de sus gastos y el 81\% dedica menos de
  una hora por semana a administrar sus finanzas.
\item
  El 68\% querría tener más control sobre sus finanzas, pero sólo si es
  fácil y no le consume tiempo.
\item
  La necesidad más mencionada es saber con exactitud cuánto gasta por
  mes y en qué (81\%).
\item
  La carga de gastos mediante WhatsApp le resultaría igual o más útil
  que la anotación manual al 81\%.
\item
  El 61\% no utiliza inteligencia artificial para sus finanzas, pero lo
  haría.
\item
  El 52\% estaría dispuesto a pagar una suscripción mensual.
\item
  16 de los 31 participantes dejaron un contacto para recibir novedades
  o colaborar con el proyecto.
\end{itemize}

El cuestionario y los resultados completos se incluyen en el Anexo A
(punto 11.2).

\textbf{Objetivos cualitativos.}

\begin{itemize}
\tightlist
\item
  Validar que las funciones que incorporan la inflación a la
  planificación financiera generan valor para los usuarios y constituyen
  un diferencial frente a la competencia.
\item
  Instalar a Quickexpense entre los jóvenes y adultos jóvenes de Rosario
  como la forma más simple de saber en qué se gasta el dinero.
\item
  Consolidar a WhatsApp como el canal principal de uso del servicio.
\end{itemize}

\textbf{Objetivos cuantitativos.} Para fijar las metas se tomaron como
referencia los siguientes indicadores de la industria:

\begin{itemize}
\tightlist
\item
  La conversión mediana a pago de las aplicaciones de suscripción con
  modelo freemium es de alrededor del 2,1\%, frente al 10,7\% de las que
  exigen el pago para acceder
  (\citeproc{ref-whitman2026revenuecat}{Whitman, 2026}).
\item
  La retención a 30 días de las aplicaciones financieras es de alrededor
  del 2\% (\citeproc{ref-mishchuk2026promodo}{Mishchuk et~al., 2026}).
\end{itemize}

{\def\LTcaptype{none} % do not increment counter
\begin{longtable}[]{@{}
  >{\raggedright\arraybackslash}p{(\linewidth - 2\tabcolsep) * \real{0.5000}}
  >{\raggedright\arraybackslash}p{(\linewidth - 2\tabcolsep) * \real{0.5000}}@{}}
\toprule\noalign{}
\begin{minipage}[b]{\linewidth}\raggedright
Indicador
\end{minipage} & \begin{minipage}[b]{\linewidth}\raggedright
Meta al cierre del primer año
\end{minipage} \\
\midrule\noalign{}
\endhead
\bottomrule\noalign{}
\endlastfoot
Usuarios registrados (acumulado) & 150 (mínimo aceptable: 100) \\
Usuarios activos mensuales & 45 (30\% de los registrados) \\
Usuarios con plan pago & 8 (5\% de los registrados) \\
Movimientos registrados mediante el chatbot & Al menos el 60\% del
total \\
\end{longtable}
}

La meta de usuarios activos supera los indicadores de referencia porque
la base inicial se construye mediante contacto directo y con usuarios
que participaron de las etapas Alpha y Beta, que suelen mostrar un
compromiso mayor que los captados mediante publicidad masiva. La
conversión a pago se fija por encima de la mediana del mercado, en línea
con la disposición a pagar declarada en la encuesta, que de todas formas
sobrestima la conversión real.

\textbf{Metas trimestrales.}

{\def\LTcaptype{none} % do not increment counter
\begin{longtable}[]{@{}
  >{\raggedright\arraybackslash}p{(\linewidth - 6\tabcolsep) * \real{0.1600}}
  >{\raggedright\arraybackslash}p{(\linewidth - 6\tabcolsep) * \real{0.3200}}
  >{\raggedright\arraybackslash}p{(\linewidth - 6\tabcolsep) * \real{0.2667}}
  >{\raggedright\arraybackslash}p{(\linewidth - 6\tabcolsep) * \real{0.2267}}@{}}
\toprule\noalign{}
\begin{minipage}[b]{\linewidth}\raggedright
Trimestre
\end{minipage} & \begin{minipage}[b]{\linewidth}\raggedright
Período
\end{minipage} & \begin{minipage}[b]{\linewidth}\raggedright
Usuarios registrados (acumulado)
\end{minipage} & \begin{minipage}[b]{\linewidth}\raggedright
Usuarios con plan pago
\end{minipage} \\
\midrule\noalign{}
\endhead
\bottomrule\noalign{}
\endlastfoot
T1 & Junio a agosto de 2027 & 60 & 2 \\
T2 & Septiembre a noviembre de 2027 & 90 & 4 \\
T3 & Diciembre de 2027 a febrero de 2028 & 120 & 6 \\
T4 & Marzo a mayo de 2028 & 150 & 8 \\
\end{longtable}
}

Antes del lanzamiento, las metas corresponden a las etapas de prueba:
hasta 20 usuarios en la versión Alpha (fines de 2026) y entre 20 y 50 en
la versión Beta (febrero de 2027). Los usuarios de la versión Beta que
continúen utilizando el servicio se contabilizan como usuarios
registrados del primer año.

\subsection*{6.2 Resultados esperados en materia de cobertura y
participación}\label{resultados-esperados-en-materia-de-cobertura-y-participaciuxf3n}
\addcontentsline{toc}{subsection}{6.2 Resultados esperados en materia de
cobertura y participación}

\textbf{Cobertura.} Durante el primer año de operación, la cobertura
geográfica se concentrará en la ciudad de Rosario. Si bien el servicio
es digital y puede utilizarse desde cualquier punto del país, las
acciones de comunicación y captación se dirigirán exclusivamente a esta
ciudad, lo que permite concentrar recursos limitados y mantener un
contacto cercano con los primeros usuarios. La expansión a otras
ciudades se evaluará en función de los resultados alcanzados.

\textbf{Participación.} El departamento Rosario cuenta con 311.156
habitantes de entre 20 y 34 años según el Censo 2022
(\citeproc{ref-correa2023censo}{Correa, 2023}), rango que constituye el
núcleo del público objetivo. La meta de 150 usuarios registrados al
cierre del primer año representa alrededor del 0,05\% de esa población.
Se trata de una participación deliberadamente acotada, coherente con una
estrategia de enfoque, con los recursos de un equipo unipersonal y con
un mercado en el que ya operan competidores locales; como referencia, el
competidor local más desarrollado declara más de 11.000 personas
usuarias en más de 20 países (\citeproc{ref-gasti2026}{Gasti, s.~f.}).

En consecuencia, el objetivo del primer año no es alcanzar una
participación significativa del mercado, sino validar la propuesta de
valor en un grupo reducido y comprometido de usuarios, que sirva como
base para el crecimiento posterior.

\capitulo{11 Anexos}

\subsection*{11.2 Anexos específicos del
proyecto}\label{anexos-especuxedficos-del-proyecto}
\addcontentsline{toc}{subsection}{11.2 Anexos específicos del proyecto}

\subsubsection*{Anexo A. Encuesta sobre hábitos de gestión de las
finanzas
personales}\label{anexo-a.-encuesta-sobre-huxe1bitos-de-gestiuxf3n-de-las-finanzas-personales}
\addcontentsline{toc}{subsubsection}{Anexo A. Encuesta sobre hábitos de
gestión de las finanzas personales}

Este anexo presenta el cuestionario y los resultados agregados de la
encuesta realizada para este proyecto
(\citeproc{ref-luraschi2026encuesta}{Luraschi, 2026}). Para resguardar
la privacidad de los participantes, no se incluyen los datos de contacto
ni las respuestas de texto libre.

\paragraph*{Ficha técnica}\label{ficha-tuxe9cnica}
\addcontentsline{toc}{paragraph}{Ficha técnica}

{\def\LTcaptype{none} % do not increment counter
\begin{longtable}[]{@{}
  >{\raggedright\arraybackslash}p{(\linewidth - 2\tabcolsep) * \real{0.2778}}
  >{\raggedright\arraybackslash}p{(\linewidth - 2\tabcolsep) * \real{0.7083}}@{}}
\toprule\noalign{}
\begin{minipage}[b]{\linewidth}\raggedright
Ítem
\end{minipage} & \begin{minipage}[b]{\linewidth}\raggedright
Detalle
\end{minipage} \\
\midrule\noalign{}
\endhead
\bottomrule\noalign{}
\endlastfoot
Período de relevamiento & 19 y 20 de septiembre de 2026 \\
Respuestas & 31 \\
Modalidad & Cuestionario en línea autoadministrado \\
Muestreo & No probabilístico, por conveniencia. Los resultados son
orientativos y no representativos de la población \\
Edad de los participantes & 19 a 27 años (media: 22,7 años; mediana: 22
años) \\
\end{longtable}
}

\paragraph*{Cuestionario}\label{cuestionario}
\addcontentsline{toc}{paragraph}{Cuestionario}

\begin{enumerate}
\def\labelenumi{\arabic{enumi}.}
\tightlist
\item
  Edad.
\item
  Mayor nivel de educación alcanzado.
\item
  ¿Trabajás?
\item
  Si trabajás como independiente o tenés tu propio negocio, por favor
  describí brevemente el rubro o actividad (respuesta abierta).
\item
  Entre las siguientes opciones, ¿cuál describe mejor tu situación?
\item
  ¿Registrás los gastos que tenés en algún lugar?
\item
  Si no registrás gastos, ¿por qué? Seleccioná todas las que apliquen.
\item
  ¿Cómo llevás el control de las cuotas de compras hechas con tarjetas
  de crédito? ¿Tenés una forma fácil de ver cómo van a impactar en el
  futuro?
\item
  Si probaste herramientas para gestionar tus finanzas, ¿cuáles usaste?
  Seleccioná todas las que apliquen.
\item
  ¿Cuánto tiempo dedicás semanalmente a revisar y administrar tus
  finanzas?
\item
  Del 1 al 10, ¿qué tan frecuentemente administrás moneda extranjera en
  tus cuentas?
\item
  ¿Tenés un presupuesto mensual?
\item
  Entre billeteras virtuales, bancos y brokers, ¿aproximadamente cuántas
  cuentas utilizás durante una semana?
\item
  ¿Tenés cuentas conjuntas o administradas por varias personas?
\item
  ¿Usás cuentas remuneradas de Mercado Pago u otra app?
\item
  ¿De qué formas invertís en el mercado? Si no lo hacés, omití la
  pregunta.
\item
  Del 1 al 10, ¿qué tanto control sentís que tenés sobre tus finanzas
  personales?
\item
  Si en algún momento atravesaste una situación financieramente
  complicada, ¿creés que hubiese sido parcial o completamente evitable
  con mejor planificación financiera?
\item
  ¿Te interesaría tener más control sobre tus finanzas?
\item
  Si en la pregunta anterior respondiste que sí, ¿cómo te gustaría
  mejorar? Seleccioná todas las que apliquen.
\item
  ¿Utilizás herramientas de inteligencia artificial con tal de mejorar
  en el seguimiento y control de tus finanzas? ¿Usarías?
\item
  Si pudieras registrar tus gastos en una aplicación enviando un mensaje
  o audio de WhatsApp, ¿qué tan útil creés que te sería?
\item
  Si una herramienta te ahorrara todo el trabajo manual de registro y te
  diera proyecciones claras de tu dinero, ¿estarías dispuesto/a a pagar
  una suscripción mensual?
\item
  Si sos emprendedor o independiente, ¿sentís que tenés problemas con
  algunos de los siguientes aspectos de tu negocio? Seleccioná todas las
  que apliquen.
\item
  Si te gustaría recibir novedades sobre el desarrollo de una
  herramienta para tus finanzas personales o las de tu emprendimiento,
  escribí alguna forma de contacto (opcional).
\item
  Si estás dispuesto/a a contestar preguntas relacionadas con finanzas
  personales para el desarrollo del proyecto, escribí alguna forma de
  contacto (opcional).
\item
  Si tenés algún comentario u observación, escribilo acá (respuesta
  abierta).
\end{enumerate}

\paragraph*{Resultados}\label{resultados}
\addcontentsline{toc}{paragraph}{Resultados}

Los porcentajes se calculan sobre el total de 31 respuestas, salvo que
se indique otra base. En las preguntas de opción múltiple, la suma puede
superar el 100\%.

\textbf{Perfil de los participantes}

{\def\LTcaptype{none} % do not increment counter
\begin{longtable}[]{@{}
  >{\raggedright\arraybackslash}p{(\linewidth - 2\tabcolsep) * \real{0.2778}}
  >{\raggedright\arraybackslash}p{(\linewidth - 2\tabcolsep) * \real{0.7083}}@{}}
\toprule\noalign{}
\begin{minipage}[b]{\linewidth}\raggedright
Pregunta
\end{minipage} & \begin{minipage}[b]{\linewidth}\raggedright
Resultados
\end{minipage} \\
\midrule\noalign{}
\endhead
\bottomrule\noalign{}
\endlastfoot
Nivel educativo & Universitario incompleto o en curso: 22 (71\%).
Técnico o terciario completo: 4 (13\%). Técnico o terciario incompleto o
en curso: 2 (6\%). Secundario completo: 2 (6\%). Universitario completo:
1 (3\%) \\
Situación laboral & No trabaja: 14 (45\%). En relación de dependencia:
10 (32\%). Independiente: 5 (16\%). Negocio o emprendimiento propio: 2
(6\%) \\
Situación de convivencia & Vive con su familia y no realiza gastos en
necesidades básicas: 16 (52\%). Vive con su familia y aporta a los
gastos generales: 7 (23\%). Vive solo/a: 5 (16\%). Vive con otra persona
y dividen gastos: 2 (6\%). Otra: 1 (3\%) \\
\end{longtable}
}

\textbf{Hábitos financieros}

{\def\LTcaptype{none} % do not increment counter
\begin{longtable}[]{@{}
  >{\raggedright\arraybackslash}p{(\linewidth - 2\tabcolsep) * \real{0.2778}}
  >{\raggedright\arraybackslash}p{(\linewidth - 2\tabcolsep) * \real{0.7083}}@{}}
\toprule\noalign{}
\begin{minipage}[b]{\linewidth}\raggedright
Pregunta
\end{minipage} & \begin{minipage}[b]{\linewidth}\raggedright
Resultados
\end{minipage} \\
\midrule\noalign{}
\endhead
\bottomrule\noalign{}
\endlastfoot
Registro de gastos & No registra ninguno: 11 (35\%). Sólo los gastos
grandes: 9 (29\%). Todo tipo de gastos: 7 (23\%). Gastos medianos y
grandes: 4 (13\%) \\
Motivos para no registrar (base: 11 que no registran) & Le parece muy
tedioso: 4. No tiene tantos gastos: 4. No lo considera necesario: 4 \\
Control de cuotas de tarjetas de crédito & No usa tarjeta de crédito: 19
(61\%). Intenta recordarlas mentalmente: 5 (16\%). Aplicación del banco:
2 (6\%). No las controla: 2 (6\%). Otras formas: 3 (10\%) \\
Herramientas utilizadas & Cuaderno o anotador: 13 (42\%). Planilla de
cálculo: 13 (42\%). ``Frascos'' de Naranja X o similares: 5 (16\%).
Cuenta separada para ahorros: 4 (13\%). Aplicaciones genéricas (Wallet,
Spendee, YNAB): 2 (6\%). Otras: 3 (10\%). Ninguna: 1 (3\%) \\
Tiempo semanal dedicado a las finanzas & Menos de 1 hora: 25 (81\%).
Entre 1 y 2 horas: 6 (19\%) \\
Frecuencia de uso de moneda extranjera (1 a 10) & Mediana: 1. Media:
2,5. 18 participantes (58\%) respondieron 1 \\
Presupuesto mensual & No: 16 (52\%). Sí, principalmente un estimado de
los gastos recurrentes: 15 (48\%) \\
Cuentas utilizadas por semana & 2 o 3: 20 (65\%). 1: 8 (26\%). 4 o 5: 2
(6\%). Ninguna: 1 (3\%) \\
Cuentas conjuntas & No: 29 (94\%). Sí: 2 (6\%) \\
Cuentas remuneradas & Sí: 12 (39\%). No está seguro/a: 12 (39\%). No: 7
(23\%) \\
Inversiones (base: 10 que invierten) & Desde la aplicación de su
billetera virtual o banco: 9. Mediante brokers como IOL: 3. Mediante
Bull Market o Cocos Capital: 3 \\
\end{longtable}
}

\textbf{Percepción y necesidades}

{\def\LTcaptype{none} % do not increment counter
\begin{longtable}[]{@{}
  >{\raggedright\arraybackslash}p{(\linewidth - 2\tabcolsep) * \real{0.2778}}
  >{\raggedright\arraybackslash}p{(\linewidth - 2\tabcolsep) * \real{0.7083}}@{}}
\toprule\noalign{}
\begin{minipage}[b]{\linewidth}\raggedright
Pregunta
\end{minipage} & \begin{minipage}[b]{\linewidth}\raggedright
Resultados
\end{minipage} \\
\midrule\noalign{}
\endhead
\bottomrule\noalign{}
\endlastfoot
Control percibido sobre sus finanzas (1 a 10) & Media: 5,45. Mediana:
5 \\
¿Una mejor planificación hubiese evitado una situación complicada? & No
le pasó: 10 (32\%). Le pasó y planificar hubiese ayudado
significativamente o la hubiese evitado: 8 (26\%). Le pasó y planificar
hubiese ayudado un poco: 7 (23\%). Le pasó y planificar no hubiese
ayudado: 6 (19\%) \\
Interés en tener más control & Sí, pero sólo si es fácil y no consume
mucho tiempo: 21 (68\%). Sí, y está dispuesto/a a hacer un esfuerzo: 10
(32\%) \\
\end{longtable}
}

{\def\LTcaptype{none} % do not increment counter
\begin{longtable}[]{@{}
  >{\raggedright\arraybackslash}p{(\linewidth - 4\tabcolsep) * \real{0.7436}}
  >{\raggedright\arraybackslash}p{(\linewidth - 4\tabcolsep) * \real{0.1667}}
  >{\raggedright\arraybackslash}p{(\linewidth - 4\tabcolsep) * \real{0.0769}}@{}}
\toprule\noalign{}
\begin{minipage}[b]{\linewidth}\raggedright
¿Cómo le gustaría mejorar?
\end{minipage} & \begin{minipage}[b]{\linewidth}\raggedright
Respuestas
\end{minipage} & \begin{minipage}[b]{\linewidth}\raggedright
\%
\end{minipage} \\
\midrule\noalign{}
\endhead
\bottomrule\noalign{}
\endlastfoot
Saber con exactitud cuánto gasta por mes y en qué & 25 & 81\% \\
Visibilizar qué inversiones son más rentables y qué oportunidades tiene
& 15 & 48\% \\
Gestionar mejor su ahorro para alcanzar ciertos objetivos & 14 & 45\% \\
Crear un presupuesto con mayor confianza y detalle & 13 & 42\% \\
Saber con exactitud cuánto gana por mes y cómo & 11 & 35\% \\
Gestionar mejor los recordatorios para pagar gastos recurrentes & 10 &
32\% \\
Recibir consejos personalizados en base a sus hábitos financieros & 10 &
32\% \\
Visibilizar el cambio de valor de sus tenencias según el tipo de cambio
y la inflación & 9 & 29\% \\
Gestionar mejor las cuotas pendientes de las tarjetas de crédito & 7 &
23\% \\
\end{longtable}
}

\textbf{Disposición hacia el producto}

{\def\LTcaptype{none} % do not increment counter
\begin{longtable}[]{@{}
  >{\raggedright\arraybackslash}p{(\linewidth - 2\tabcolsep) * \real{0.2778}}
  >{\raggedright\arraybackslash}p{(\linewidth - 2\tabcolsep) * \real{0.7083}}@{}}
\toprule\noalign{}
\begin{minipage}[b]{\linewidth}\raggedright
Pregunta
\end{minipage} & \begin{minipage}[b]{\linewidth}\raggedright
Resultados
\end{minipage} \\
\midrule\noalign{}
\endhead
\bottomrule\noalign{}
\endlastfoot
Uso de inteligencia artificial para sus finanzas & No usó, pero usaría:
19 (61\%). No usó y no usaría: 7 (23\%). Usa y seguiría usando: 3
(10\%). Usó o usa, pero preferiría dejar de usarla: 2 (6\%) \\
Utilidad de registrar gastos por mensaje o audio de WhatsApp & Más útil
que anotar el gasto a mano: 15 (48\%). Igual de útil: 10 (32\%). Muy
poco útil: 6 (19\%) \\
Disposición a pagar una suscripción mensual & Sí: 16 (52\%). No: 15
(48\%) \\
\end{longtable}
}

\textbf{Emprendedores e independientes} (base: 8 respuestas): cobranza
de deudas, 7; registro de ventas, 5; gestión y ajuste de precios, 3;
registro de compras, 3.

\subsection*{11.3 Bibliografía
Consultada}\label{bibliografuxeda-consultada}
\addcontentsline{toc}{subsection}{11.3 Bibliografía Consultada}

\protect\phantomsection\label{refs}
\begin{CSLReferences}{1}{0}
\bibitem[\citeproctext]{ref-aimoney}
AI Money. (s.~f.). \emph{La mejor app de finanzas y gastos con IA,
gratis}. Recuperado 5 de octubre de 2026, de
\url{https://www.ai-money.app/}

\bibitem[\citeproctext]{ref-bcracaf2018}
Banco Central de la República Argentina, \& CAF - Banco de Desarrollo de
América Latina. (2018). \emph{Encuesta de medición de capacidades
financieras en Argentina: 2017}.
\url{https://scioteca.caf.com/handle/123456789/1341}

\bibitem[\citeproctext]{ref-chanchito}
Chanchito. (s.~f.). \emph{Las mejores apps para controlar gastos en
2026}. Recuperado 5 de octubre de 2026, de
\url{https://www.chanchito.app/blog/app-para-controlar-gastos}

\bibitem[\citeproctext]{ref-correa2023censo}
Correa, G. (2023, noviembre 25). Cuántos somos, cómo estamos: deudas de
los 40 años de democracia en el departamento Rosario. \emph{El
Ciudadano}.
\url{https://elciudadanoweb.com/cuantos-somos-como-estamos-deudas-de-los-40-anos-de-democracia-en-el-departamento-rosario/}

\bibitem[\citeproctext]{ref-decreto353}
Decreto 353/2025, Boletín Oficial de la República Argentina (2025).
\url{https://www.boletinoficial.gob.ar/detalleAviso/primera/325767/20250523}

\bibitem[\citeproctext]{ref-mdz2026emae}
Día D para la economía: el Indec revela el EMAE de julio y se define si
Argentina entra en recesión. (2026, septiembre 24). \emph{MDZ}.
\url{https://www.mdzol.com/dinero/dia-d-la-economia-el-indec-revela-el-emae-julio-y-se-define-si-argentina-entra-recesion-n1613691}

\bibitem[\citeproctext]{ref-dpa1751}
Digital Policy Alert. (s.~f.). \emph{Data protection regulation in bill
providing regulatory framework for protection of personal data and
privacy (No.~1751-D-2026)}. Recuperado 6 de octubre de 2026, de
\url{https://digitalpolicyalert.org/change/19257-data-protection-regulation-in-bill-providing-regulatory-framework-for-protection-of-personal-data-and-privacy-no-1751-d-2026}

\bibitem[\citeproctext]{ref-monthly2026}
Equipo Monthly. (2026, junio 1). Las mejores apps de finanzas personales
para Android en Argentina. \emph{Monthly}.
\url{https://monthly.com.ar/blog/mejores-apps-finanzas-personales-android-argentina/}

\bibitem[\citeproctext]{ref-piggy2026}
Equipo Piggy. (2026, septiembre 17). \emph{¿Cuál es la mejor app de
finanzas personales en Argentina en 2026?} Piggy.
\url{https://piggy.com.ar/guias/mejor-app-finanzas-personales-argentina/}

\bibitem[\citeproctext]{ref-filgueira2025}
Filgueira, R. (2025, octubre 27). Not all chatbots are banned:
WhatsApp's 2026 AI policy explained. \emph{respond.io}.
\url{https://respond.io/blog/whatsapp-general-purpose-chatbots-ban}

\bibitem[\citeproctext]{ref-fiskil2026}
Fiskil. (s.~f.). \emph{Open Banking Argentina: guía completa BCRA para
fintechs}. Recuperado 6 de octubre de 2026, de
\url{https://www.fiskil.com/es/open-finance/solutions/argentina/overview}

\bibitem[\citeproctext]{ref-focusfolio2026}
Focus·Folio. (2026). \emph{Las mejores apps de finanzas personales en
Argentina (2026)}.
\url{https://focusfolio.com.ar/mejores-apps-finanzas-personales-argentina}

\bibitem[\citeproctext]{ref-gasalla2026dolar}
Gasalla, J. (2026, octubre 5). El dólar operó sin cambios y quedó a más
de 400 pesos del techo de las bandas cambiarias. \emph{Infobae}.
\url{https://www.infobae.com/economia/2026/10/05/el-dolar-opero-sin-cambios-y-quedo-a-mas-de-400-pesos-del-techo-de-las-bandas-cambiarias/}

\bibitem[\citeproctext]{ref-gasti2026}
Gasti. (s.~f.). \emph{Gasti --- Finanzas personales, más allá de
trackear gastos}. Recuperado 6 de octubre de 2026, de
\url{https://gasti.pro/}

\bibitem[\citeproctext]{ref-gimenez2026bandas}
Giménez, J. (2026, enero 2). Dólar y nuevo esquema cambiario: cómo
quedan las bandas de flotación a partir del 2 de enero de 2026.
\emph{Chequeado}.
\url{https://chequeado.com/el-explicador/bandas-de-flotacion-del-dolar-que-cambia-desde-2026-y-como-funcionaron-hasta-ahora/}

\bibitem[\citeproctext]{ref-indec2025tic}
Instituto Nacional de Estadística y Censos. (2025). \emph{Acceso y uso
de tecnologías de la información y la comunicación. EPH. Cuarto
trimestre de 2024} (No. 113; Informes técnicos, Vol. 9).
\url{https://www.indec.gob.ar/uploads/informesdeprensa/mautic_05_25FD0D0C4A71.pdf}

\bibitem[\citeproctext]{ref-indec2026emae}
Instituto Nacional de Estadística y Censos. (2026a). \emph{Estimador
mensual de actividad económica. Julio de 2026} (Informes técnicos).
\url{https://www.indec.gob.ar/uploads/informesdeprensa/emae_09_26FE154D8FFE.pdf}

\bibitem[\citeproctext]{ref-indec2026ipc}
Instituto Nacional de Estadística y Censos. (2026b). \emph{Índice de
precios al consumidor (IPC). Agosto de 2026} (Informes técnicos).
\url{https://www.indec.gob.ar/uploads/informesdeprensa/ipc_09_26A1BE2DC4CD.pdf}

\bibitem[\citeproctext]{ref-indec2026pib}
Instituto Nacional de Estadística y Censos. (2026c). \emph{Informe de
avance del nivel de actividad. Segundo trimestre de 2026} (Informes
técnicos).
\url{https://www.indec.gob.ar/uploads/informesdeprensa/pib_09_26065D6A3612.pdf}

\bibitem[\citeproctext]{ref-indec2026eph}
Instituto Nacional de Estadística y Censos. (2026d). \emph{Mercado de
trabajo. Tasas e indicadores socioeconómicos (EPH). Segundo trimestre de
2026} (Informes técnicos).
\url{https://www.indec.gob.ar/uploads/informesdeprensa/mercado_trabajo_eph_2trim26433FCBC5A8.pdf}

\bibitem[\citeproctext]{ref-kemp2025digital}
Kemp, S. (2025). \emph{Digital 2026: Argentina}. DataReportal.
\url{https://datareportal.com/reports/digital-2026-argentina}

\bibitem[\citeproctext]{ref-kotler2016}
Kotler, P., \& Keller, K. L. (2016). \emph{Marketing management} (15.ª
ed.). Pearson.

\bibitem[\citeproctext]{ref-lepiscopo2026iva}
L'Episcopo, I. (2026, octubre). \emph{IVA RG 4240: qué es y cómo se
cobra en Netflix y Spotify 2026}. nacho.com.ar.
\url{https://nacho.com.ar/iva-netflix-cobraran-el-impuesto-a-los-usuarios-de-servicios-digitales/}

\bibitem[\citeproctext]{ref-ley25326}
Ley 25.326 de Protección de los Datos Personales, Boletín Oficial de la
República Argentina (2000).
\url{https://servicios.infoleg.gob.ar/infolegInternet/anexos/60000-64999/64790/texact.htm}

\bibitem[\citeproctext]{ref-luraschi2026encuesta}
Luraschi, M. (2026). \emph{Encuesta sobre hábitos de gestión de las
finanzas personales} {[}Conjunto de datos no publicado{]}.

\bibitem[\citeproctext]{ref-lanacion2026pbi}
Manzoni, C. (2026, marzo 20). La economía creció 4,4\% en 2025,
impulsada por los bancos, el agro y la minería. \emph{La Nación}.
\url{https://www.lanacion.com.ar/arquitectura/la-economia-crecio-44-en-2025-con-sectores-ganadores-y-perdedores-nid20032026/}

\bibitem[\citeproctext]{ref-marval2026datos}
Marval O'Farrell Mairal. (2026, agosto 27). \emph{Nuevo proyecto de ley
de protección de datos}.
\url{https://www.marval.com/Publicacion/nuevo-proyecto-de-ley-de-proteccion-de-datos-17691}

\bibitem[\citeproctext]{ref-mishchuk2026promodo}
Mishchuk, Y., Loshak, V., \& Afanasieva, Y. (2026, julio 31). Mobile
marketing benchmarks: Know what good performance looks like in 2026.
\emph{Promodo}.
\url{https://www.promodo.com/blog/mobile-marketing-benchmarks}

\bibitem[\citeproctext]{ref-porter1980}
Porter, M. E. (1980). \emph{Competitive strategy: Techniques for
analyzing industries and competitors}. Free Press.

\bibitem[\citeproctext]{ref-cronica2026ia}
Presentan un proyecto para convertir a Argentina en un polo global de
IA. (2026, junio 14). \emph{Diario Crónica}.
\url{https://www.diariocronica.com.ar/noticias/2026/06/14/149011-presentan-un-proyecto-para-convertir-a-argentina-en-un-polo-global-de-ia}

\bibitem[\citeproctext]{ref-quewallet}
QueWallet. (s.~f.). \emph{QueWallet - Control de gastos por WhatsApp}.
Recuperado 5 de octubre de 2026, de \url{https://www.quewallet.com/}

\bibitem[\citeproctext]{ref-statista2025redes}
Statista. (2025). \emph{Most popular social media platforms in Argentina
as of 2nd quarter 2025, by usage reach} {[}Gráfico{]}.
\url{https://www.statista.com/statistics/284401/argentina-social-network-penetration/}

\bibitem[\citeproctext]{ref-szlak2025}
Szlak, G. (2025, junio 3). Novedades legislativas en Argentina sobre
protección de datos personales e inteligencia artificial. \emph{IAPP}.
\url{https://iapp.org/news/a/novedades-legislativas-en-argentina-sobre-protecci-n-de-datos-personales-e-inteligencia-artificial}

\bibitem[\citeproctext]{ref-utdt2025ia}
Universidad Torcuato Di Tella. (2025, noviembre 18). \emph{Cómo y para
qué usamos la inteligencia artificial en Argentina, según un estudio
académico}.
\url{https://www.utdt.edu/ver_nota_prensa.php?id_nota_prensa=23080}

\bibitem[\citeproctext]{ref-whatsapp-terms}
WhatsApp. (s.~f.). \emph{WhatsApp Business Solution Terms}. Recuperado 6
de octubre de 2026, de
\url{https://www.whatsapp.com/legal/business-solution-terms}

\bibitem[\citeproctext]{ref-whitman2026revenuecat}
Whitman, L. (2026, marzo 19). The state of subscription apps in 10
minutes: Lessons, trends, and benchmarks for 2026. \emph{RevenueCat}.
\url{https://www.revenuecat.com/blog/growth/subscription-app-trends-benchmarks-2026}

\bibitem[\citeproctext]{ref-ynab-pricing}
YNAB. (s.~f.-a). \emph{Pricing}. Recuperado 6 de octubre de 2026, de
\url{https://www.ynab.com/pricing}

\bibitem[\citeproctext]{ref-ynab-monedas}
YNAB. (s.~f.-b). \emph{Using multiple currencies in YNAB: A guide}.
Recuperado 7 de octubre de 2026, de
\url{https://support.ynab.com/en_us/using-multiple-currencies-in-ynab-a-guide-SyBF6PHno}

\end{CSLReferences}
